% Transcription — fonds Alexandre Grothendieck, shelfmark 162-5, batch 1 (pages 1-20).
% Established from the facsimile archives/batches/162-5/batch-01.pdf (a mirror of
% https://grothendieck.umontpellier.fr/162-5.pdf), per the skill
% transcribe-grothendieck. The batch file carries no cover sheet: PDF page N is
% archivists' page N, as the pencilled numbers bottom-left confirm.
%
% Pass: Opus 5.5 (claude-opus-5-5), 2026-09-22 — first pass, unchecked against the pages by a human.
% Revised: Opus 5.5 (claude-opus-5-5), 2026-09-22 — shared ink layer read against copy B (pp. 32–48).
%
% What the batch is.
%   1      The folder's cover: « Le Tapis de Quillen / notes d'Alexandre
%          (1968) », in pencil, in a hand that is probably not his (he would
%          hardly write « d'Alexandre »). No \page; its words head the file
%          and a \note says where they come from.
%   2-17   A typescript of sixteen leaves, paginated 1 to 16 at the head
%          (typescript page N = archivists' page N+1), headed « Tapis de
%          Quillen » and dated « notes du 10.9.68 ». It is his: first person
%          throughout (« mon interprétation chérie », « ma lettre à Verdier »,
%          « Je pense que ce résultat (facile) doit pouvoir se généraliser »),
%          reporting what Quillen told him and what to ask Quillen next.
%   18, 20 Blank separator sheets: absent.
%   19     A typed list, « Publications of Daniel Quillen », items 1-11 (1966
%          to "to appear"), with nothing in his hand: absent, like any page
%          that is not the mathematics in hand. It is not his and carries no
%          annotation; the tiré à part the inventory mentions is presumably in
%          a later batch.
%
% What copy this is. Folder 162-5 holds two witnesses of this typescript.
% Copy B (pp. 32-40 in batch 2, pp. 41, 43-48 in batch 3) is the typed top
% copy annotated in blue ink. These pages (copy A) are a black photocopy of
% copy B taken after the blue ink was added: same line breaks, same typewriter
% x-outs, the right edge clipped on most leaves. Copy A typescript page N =
% archivists' page N+1; copy B typescript page N = page 31+N for N <= 9, page
% 41 for N = 10, page 32+N for N = 11-16. Copy A's hand layer is therefore two
% layers: (1) B's ink, photocopied — the date, the framed note of p. 8, the
% interlinear of p. 8 and its margin, the margin answer of p. 9, the
% « Attention » of p. 4, the handwritten paragraph of p. 11, the insertions and
% the margin note of pp. 12-17, the closing paragraph of p. 17; (2) marks
% only copy A carries, absent from B: « Si ! » with its brace (p. 4), « t une »
% over « définie en » and « cond. » over struck « colonne » (p. 9),
% « (compactes ?) » (p. 12), « NB H*(X) (= H*(X, F_2)) » and « i.e.
% l'enveloppe "karoubienne" » (p. 14), « = karoubienne + additive », « en »
% over « considérant », and Λ inked into the two blanks « sur Λ. » and
% « puisque Λ n'est pas un corps » (p. 15), « él. des » (p. 16), the t of
% « Quillent » struck and « cf exposé Karoubi à Bourbaki ! » (p. 17).
% Folder 111 is a third reproduction of the same annotated state and carries
% the same A-only marks.
%
% Revision of 2026-09-22. The first pass read the photocopy alone and misread
% part of the shared layer. Every handwritten element was then read against
% copy B (pp. 32-48) and batches 2-3: where copy A's own image is consistent
% with B's clearer ink, the reading follows B; where A's image cannot show it
% (edge clipped, photocopy pale), \uncertain / \ill stay and nothing is
% imported. Apparatus conventions follow batches 2-3 for the same ink: \add
% for ink inserted into the typed line, \struck for typed words struck by
% hand, \marginal for detached notes; typed interlinear insertions (« la
% catégorie homotopique », « d'anneaux gradués », « stable », « fixe », « co »,
% « i.e. commutative (car. 2 !) », « additifs », « S(X) ») are typed text and
% set plain; typewriter x-outs are not reproduced. The hand's small curled
% reference signs in the left margin are not reproduced.
%
% Mathematics, in three numbered parts of the typescript.
%   I   (pp. 2-6) categories vs semi-simplicial sets: the nerve S : Cat →
%       Ssimpl, the category of simplices T(X) = Δ/X, Quillen's result that
%       S and T induce quasi-inverse equivalences of the homotopy categories,
%       the canonical homotopism C ≃ C°; local systems and lim^(i), the topos
%       C~; the equivalence of D^b_lc induced by a homotopism, its hoped-for
%       generalisation to morphisms of « locally homotopically trivial »
%       topoi (with his own note that étale topologies fail the condition),
%       and the equivalent conditions a), b), a') on p. 5.
%   II  (pp. 6-11) n-categories, n-fold categories, Gr-categories: bicategories
%       as internal categories in Cat ↔ bi-semisimplicial sets; 2-categories;
%       (Cat)-groups, Dold–Kan, and Picard categories ↔ complexes of length 1;
%       strict abelian 2-Cat-groups ↔ complexes of length 2; the general hope
%       that abelian groups in n-Cat ↔ complexes of length n; K_n(C) for a
%       monoidal category as π_n of a universal Gr-n-category.
%   III (pp. 12-17) « point de vue motivique en théorie du cobordisme »:
%       unoriented bordism B(X), the universal pair (F_•, F^•) with base
%       change for transversal squares, the correspondence category B,
%       B_•(X) ≃ H^•(X) ⊗ Λ, a « Tate object » Λ(-1) and the completion B~,
%       the analogy with motives modulo rational equivalence; then
%       S(X) ⊂ H^•(X) (classes of submanifolds), Thom's theorem as
%       invariance under the Steenrod group scheme G over F_2 read as a
%       « motivic Galois group », non-reductive, pro-unipotent (per Quillen);
%       oriented variants via H → BO(∞).
%
% Notation fixed for the batch. The typescript writes « . » and « ' » raised
% for lower and upper bullets (F., F', f., f', B.(X), H'(X)): set as
% F_\bullet, F^\bullet, B_\bullet(X), H^\bullet(X). C° (opposite category) is
% C^\circ. The underlined bold F_2 and Q are \underline{\mathbb{F}}_2,
% \underline{\mathbb{Q}}; Hom underlined is \underline{\mathrm{Hom}}; ⊗ is
% typed as an overstruck O and given as \otimes; « lim » with a left arrow
% under it is \varprojlim. His hand rewrites the indices of C_0, C_1, X_ij,
% typed as superscripts, as subscripts; the correction is followed.
%
% Readings left open: the oblique left-margin blocks of pp. 8 and 17, whose
% photocopy is pale and clipped at the left edge; the condition on the
% generators x_i (p. 14); the struck word and the exponent before H^1(S^1)
% (p. 16); the word after « bien par » in the « Attention » of p. 4. The
% typescript's right edge is clipped on most leaves; missing word-ends are
% restored with \add.
%
% The argument does not stop at p. 17: the typescript breaks off on its last
% line and copy A closes with its own « cf exposé Karoubi à Bourbaki » ;
% nothing in pp. 18-20 carries it on. Whatever follows in the folder (batch 2)
% begins after the list of Quillen's publications.
\documentclass[11pt,a4paper]{article}
% Le préambule commun à toutes les transcriptions du fonds Grothendieck.
%
% Il tient en une idée : **l'incertitude se compose, elle ne se résout pas.**
% Une transcription qui lisse un mot douteux a détruit la seule chose pour
% laquelle on la faisait — un lecteur doit pouvoir savoir, à chaque endroit,
% ce qui était sur le feuillet et ce qui a été ajouté. Les six macros qui
% suivent sont donc l'appareil critique, et non de la décoration.
%
% Les mêmes commandes sont reconnues par `scripts/render.mjs`, qui produit la
% vue de lecture du site. Une macro ajoutée ici doit l'être aussi là-bas,
% faute de quoi le rendu échouera bruyamment — ce qui est voulu : un rendu
% silencieusement faux est pire qu'un rendu absent.

\ProvidesPackage{grothendieck}[2026/08/08 Transcriptions du fonds Grothendieck]

\usepackage{amsmath,amssymb,amsthm}
\usepackage{mathrsfs}
\usepackage{stmaryrd}
% Les diagrammes commutatifs sont partout dans ce fonds : c'est la langue
% même de la géométrie algébrique de Grothendieck, et une transcription qui
% les rendrait en prose ne transcrirait rien.
\usepackage{tikz-cd}
% Français par défaut — c'est la langue des feuillets — mais l'anglais est
% chargé aussi : la traduction et le résumé sont des documents anglais, et la
% typographie française (espaces avant les deux-points) y serait une faute.
% Ces documents basculent par \selectlanguage{english} en tête de corps.
\usepackage[english,french]{babel}
\usepackage{fontspec}
\usepackage{geometry}
\geometry{a4paper,margin=25mm,marginparwidth=28mm}
\usepackage{marginnote}
\usepackage{xcolor}
% ulem, not soul. `soul`'s \st and \ul are built on a character-by-character
% scan that XeLaTeX's font machinery breaks: the document fails with an
% undefined control sequence rather than degrading. ulem does the same two jobs
% and is Unicode-engine safe. `normalem` stops it hijacking \emph, which the
% transcriptions use for Grothendieck's own emphasis.
\usepackage[normalem]{ulem}
\usepackage{hyperref}
\hypersetup{colorlinks=true,linkcolor=black,urlcolor=black,pdfborder={0 0 0}}

% --- Métadonnées du lot -----------------------------------------------------
% Reprises telles quelles par le script de rendu, qui les affiche en tête de
% la vue de lecture. Elles fixent ce que le fichier prétend couvrir : sans
% elles, une transcription flotte sans point d'attache dans un fonds de seize
% mille feuillets.

\newcommand{\folder}[1]{\gdef\grothFolder{#1}}
\newcommand{\batch}[1]{\gdef\grothBatch{#1}}
\newcommand{\leaves}[2]{\gdef\grothFirst{#1}\gdef\grothLast{#2}}
\newcommand{\foldertitle}[1]{\gdef\grothFoldertitle{#1}}
\gdef\grothFolder{?}\gdef\grothBatch{?}\gdef\grothFirst{?}\gdef\grothLast{?}\gdef\grothFoldertitle{}

% --- Le résumé --------------------------------------------------------------
%
% Il ouvre la lecture modernisée, sous le titre « Résumé », et il est composé
% en petit. Ce n'est pas un ornement : c'est le seul passage du document qui
% ne soit pas des mathématiques, et un lecteur doit pouvoir le reconnaître
% avant de l'avoir lu — puis le sauter s'il connaît déjà le sujet, ou s'y
% arrêter s'il ne le connaît pas. Le filet qui le suit marque la reprise du
% corps.

\newenvironment{resume}
  {\small\setlength{\parskip}{0.45em}}
  {\par\medskip\hrule\medskip}

% --- Mots-clés ---------------------------------------------------------------
%
% En anglais, en clôture du résumé de la lecture modernisée : le vocabulaire
% moderne sous lequel chercher ce que les feuillets construisent. Le manifeste
% du site (`npm run manifest`) les extrait pour étiqueter la cote entière —
% cette ligne est la source unique des tags, il n'y a pas de fichier à côté.
\newcommand{\keywords}[1]{\par\smallskip\noindent
  {\footnotesize\textsc{Keywords} --- \textit{#1}}\par}

% --- L'appareil critique ----------------------------------------------------

% Le numéro de feuillet, en marge. C'est l'ancre qui permet de régler un
% désaccord sur une formule en regardant *un* feuillet plutôt qu'une chemise
% de six cents. Le site s'en sert aussi pour tourner les pages du fac-similé
% quand on fait défiler la transcription.
\newcommand{\leaf}[1]{%
  \marginnote{\footnotesize\color{gray!70}\textsf{#1}}%
  \label{leaf:#1}\ignorespaces}

% Illisible : on ne devine pas. Un mot inventé qui se lit comme les autres est
% le pire résultat possible.
\newcommand{\ill}{\textcolor{red!60!black}{[\,\ldots\,]}}

% Lecture douteuse : le mot est donné, et signalé comme incertain.
\newcommand{\uncertain}[1]{\uline{#1}}

% Ajout de l'éditeur — une lettre manquante, un mot sous-entendu.
\newcommand{\add}[1]{\textcolor{blue!50!black}{[#1]}}

% Biffé par Grothendieck lui-même. Ce qu'il a rayé fait partie du document :
% une rature dit souvent où la pensée a changé de direction.
\newcommand{\struck}[1]{\sout{#1}}

% Note du transcripteur, sur l'état matériel du feuillet.
\newcommand{\note}[1]{\par\smallskip{\small\color{gray!60!black}\textit{[#1]}}\par\smallskip}

% Note marginale de Grothendieck, distincte du corps du texte.
\newcommand{\marginal}[1]{\par\smallskip\noindent\hspace{1em}%
  {\small\color{gray!60!black}#1}\par\smallskip}

% --- Titre ------------------------------------------------------------------

\AtBeginDocument{%
  \begin{center}
    {\large\bfseries Fonds Alexandre Grothendieck --- cote n\textsuperscript{o}~\grothFolder}\\[2pt]
    {\itshape\grothFoldertitle}\\[4pt]
    {\small feuillets \grothFirst--\grothLast\ (lot \grothBatch)}\\[2pt]
    {\footnotesize\color{gray!60!black}Transcription établie d'après le fac-similé numérisé
     par l'Université de Montpellier.\\
     Les lectures incertaines sont soulignées, les illisibles notées [\,\ldots\,],
     les ajouts entre crochets.}
  \end{center}
  \vspace{1em}}

\endinput


\folder{162-5}
\batch{1}
\pages{1}{20}
\dating{1968-[à partir de 1970] — le groupe « Dossiers rassemblés par Grothendieck sur différentes thématiques » (161-1 à 162-6) est daté [après 1961-vers 1977]}
\watermark{Édition de démonstration}
\foldertitle{Tapis de Quillen : tapuscrit et copies de tapuscrits (1968, s.d.), notes manuscrites (1968), tiré à part (1968).}

\begin{document}

\section*{Le Tapis de Quillen}

\note{Ces mots, avec au-dessous « notes d'Alexandre (1968) », sont écrits au
crayon sur la chemise (page 1), qui ne porte rien d'autre ; la main n'est pas
sûrement la sienne. Suit un tapuscrit paginé 1 à 16 en tête de page, que
l'archiviste numérote 2 à 17 : la page $N$ du tapuscrit est la page $N+1$
ci-dessous. Le texte parle à la première personne (« mon interprétation
chérie », « le rôle universel que j'ai envisagé plus haut ») et rapporte ce que
Quillen lui a exposé ; il est annoté à la main.}

\page{2}
\marginal{notes du 10.9.68}

Tapis de Quillen.

\subsection*{I \quad Relations entre catégories et ensembles semi-simpliciaux.}

A toute catégorie $C$, on associe un ensemble semi-simplicial $S(C)$,
trouvant ainsi un foncteur pleinement fidèle
\[
S\,;\ (\mathrm{Cat}) \longrightarrow (\mathrm{Ssimpl}) \quad .
\]
Les systèmes locaux d'ensemble sur $SC$ correspondent aux foncteurs sur $C$
qui transforment toute flèche en isomorphisme (i.e. qui se factorisent par le
groupoïde associé à $C$). Les $H^{i}$ sur $SC$ d'un tel système local ($H^{0}$
pour ens, $H^{1}$ pour groupes, $H^{i}$ quelconques pour groupes abéliens)
s'interprètent en termes des foncteurs $\varprojlim^{(i)}$ dérivés de
$\varprojlim$, ou si on préfère, des $H^{i}($ du \emph{topos} $C$\,\add{)}. On voit
ainsi à quelle condition un foncteur $C \to C'$ induit un homotopisme
$SC \to SC'$ : en vertu du critère cohomologique de Artin-Mazur, il f\add{au}t
et s\add{uffit} que pour tout système de coéfficients $F'$ sur $C'$,
l'homomorphisme naturel $\varprojlim^{(i)}_{C'} F' \to \varprojlim^{(i)}_{C} F$
soit un isomorphisme (pour les $i$ pour lesquels cela a un sens).

A $C$ on peut asso\add{ci}er le topos $\widetilde{C}$, qui varie de façon
\emph{covariante} avec $C$ (NB le foncteur $C \mapsto \widetilde{C}$ n'a plus
rien de pleinement fidèle, semble-t-il ??) Les systèmes de coéfficients
ensemblistes sur $C$ ( = les foncteurs $C^{\circ} \to \mathrm{Ens}$
transformant isomorphismes en isomorphismes) correspondent aux faisceaux
localement constants i.e. les objets localement constants de $\widetilde{C}$,
définis intrinsèquement en termes de $\widetilde{C}$. Ainsi, le fait pour un
foncteur $F : C \to C'$ d'induire un homotopisme $S(C) \to S(C')$ ne dépend que
du morphisme de topos $\widetilde{F} : \widetilde{C} \to \widetilde{C}'$
induit, et signifie que pour tout faisceau localement constant $F'$ sur $C'$
i.e. sur $\widetilde{C}'$, les applications induites
$H^{i}(\widetilde{C}', F') \to H^{i}(\widetilde{C}, \widetilde{F}^{*}(F'))$
sont des isomorphismes (pour les $i$ pour lesquels cela a un sens).
\note{Le
second membre de la flèche $\varprojlim^{(i)}_{C'} F' \to \varprojlim^{(i)}_{C} F$
est tapé $F$ et non l'image inverse de $F'$ ; on le laisse tel quel.}

\page{3}
On a aussi un foncteur évident
\[
T : (\mathrm{Ssimpl}) \longrightarrow (\mathrm{Cat}) ,
\]
en associant à tout ensemble semi-simplicial $X$ la catégorie
$T(X) = \Delta_{/X}$ des simplexes sur $X$, dont l'ensemble des objets est la
réunion disjointe des $X_{n}$ … (c'est une catégorie fibrée sur la
catégorie $\Delta$ des simplexes types, à fibres les catégories discrètes
définies par les $X_{n}$). Ceci posé Quillen prouve que pour tout $X$, $ST(X)$
est isomorphe canoniquement à $X$ dans la catégorie homotopique construite avec
$(\mathrm{Ssimpl})$, et que pour toute $C$, la catégorie $TS(C)$ est
canoniquement « homotopiquement équivalente à $C$ » i.e canoniquement
isomorphe à $C$ dans la catégorie quotient de $(\mathrm{Cat})$ obtenue en
inversant le\add{s} foncteurs qui sont des homotopismes. Ces isomorphismes
son\add{t} fonctoriels en $X$. Il en résulte formellement qu'un morphisme
$f : X \to Y$ dans $(\mathrm{Ssimpl})$ est un homotopisme si et seulement si il
en est ainsi de $T(f) : T(X) \to T(Y)$, d'où des foncteurs
$S' : (\mathrm{Cat})' \to (\mathrm{Ssimpl})'$ et
$T' : (\mathrm{Ssimpl})' \to (\mathrm{Cat})'$ entre les catégories
« homotopiques », construites avec $(\mathrm{Cat})$ resp.\ $(\mathrm{Ssimpl})$,
qui sont quasi-inverses l'un de l'autre.

De plus, Quillen construit un isomorphisme canonique et fonctoriel dans
$(\mathrm{Cat})'$ entre $C$ et la catégorie opposée $C^{\circ}$, ou ce qui
revient au même, un isomorphisme canonique et fonctoriel dans
$(\mathrm{Ssimpl})'$ entre $S(C)$ et $S(C^{\circ})$. La définition est telle
que le foncteur induit sur les systèmes locaux sur $C$ transforme le foncteur
contravariant $F$ sur $C$, transformant toute flèche en flèche inversible, en
le foncteur covariant (i.e. contravariant sur $C^{\circ}$) ayant mêmes valeurs
sur les objets, et obtenu sur les flèches en remplaçant $F(u)$ par
$F(u)^{-1}$ ; en d'autres termes, l'effet de l'homotopisme de Quillen sur les
groupoïdes fondamentaux est l'isomorphisme évident entre les groupoïdes
fondamentaux de $C$ et de $C^{\circ}$, compte tenu que le deuxième est
l'opposé du

\page{4}
premier. Comme application, Quillen obtient une interprétation faisceautique
de la cohomologie d'un ensemble semi-simplicial à coéfficients dans un système
local covariant $F$ (défini classiquement par le complexe cosimplicial des
$C^{n}(F) = \coprod_{x \in X_{n}} F(x)$) : on considère le système local
contravariant défini par $F$, on l'interprète comme un faisceau sur $T(X)$ i.e.
objet de $(\mathrm{Ssimpl})_{/X}$, et on prend sa cohomologie.
\marginal{$\prod$ ?}
\note{en marge de gauche, à hauteur de $C^{n}(F) = \coprod F(x)$ : la marge
demande un produit à la place du signe tapé $\coprod$. « covariant » est ce qui
reste d'une frappe « contravariant » dont la machine a biffé la syllabe
« tra ».}
--- Cependant, quand $F$ est un système de coéfficients covariant pas
nécessairement local, on n'a tjrs pas d'interprétation de ses groupes de
cohomologie classiques en termes faisceautiques ; ni, lorsque $F$ est
contravariant, de son homologie, ou inversement de sa cohomologie faisceautique
en termes classiques.
\marginal{Si !}
\note{en marge de gauche, au milieu d'une accolade qui embrasse toute la
phrase « Cependant, quand $F$ … en termes classiques » ; cette réponse est
propre à cette copie, l'exemplaire annoté à l'encre (p.~34) ne la porte pas.}

A propos de la notion de foncteur qui est un homotopisme. Quillen montre qu'un
tel foncteur $F : C \to C'$ induit une équivalence entre la sous catégorie
triangulée $D^{b}_{\mathrm{lc}}(C')$ de la catégorie dérivée bornée de celle
des faisceaux abéliens sur $C'$, dont les faisceaux de cohomologie sont des
systèmes locaux, et la catégorie analogue pour $C$ ; et réciproquement. On
peut dans cet énoncé introduire aussi n'importe quel anneau de base (à
conditi\add{on} de le supposer $\neq 0$ dans le cas de la réciproque) ; la
partie dir\add{ecte} vaut aussi avec un anneau de coefficients pas néc.\ 
constant, mais consta\add{nt} tordu. Je pense que ce résultat (facile) doit
pouvoir se généraliser ai\add{nsi}.

Soit $f : X \to X'$ un morphisme de topos qui soit tel que pour tout faisceau
loc.\ const.\ sur $X'$, $f$ induise un isomorphisme sur les cohomolog\add{ies}
(avec cas non commutatif inclus). Supposons que \struck{\ill{}} \add{$X$ et
$X'$} soit \emph{locale\add{ment} homotopiquement trivial}, i.e. que pour tout
entier $n$ \add{$\geq 1$}, tout objet $U$ \struck{\ill{} $X'$} ait un
recouvrement par des $U_{i} \to U$, tels que a) tout système local sur
\add{$U$} devient constant sur $U_{i}$, et toute section sur $U$ devient
constante sur $U$\add{$_{i}$} et b) pour tout groupe abélien $G$, les
$H^{j}(U, G) \to H^{j}(U_{i}, G)$ son\add{t} nuls pour $1 \leq j \leq n$. Alors
le foncteur $D^{b}_{\mathrm{lc}}(X') \to D^{b}_{\mathrm{lc}}(X)$ induit par
$f$ est
\note{la définition, de « Supposons que » à « $1 \leq j \leq n$ », est tenue
en marge par un double trait vertical et un renvoi à la note qui suit, écrite
à la main au bas de la page ; le bord droit du tapuscrit est rogné, d'où les
fins de mots restituées. « $X$ et $X'$ » est écrit à la main au-dessus d'un
mot tapé qu'il biffe ; après « tout objet $U$ », la fin de ligne est rognée,
et la ligne suivante s'ouvre sur un $X'$ biffé à la main.}
\marginal{Attention, cette condition n'est \uncertain{typiquement}
\emph{pas} satisfaite par les schémas avec leur topologie étale (mais bien
par \struck{\ill{}} \uncertain{variétés} loc.\ noeth.\ avec top.\ Zariski).}


\page{5}
une équivalence). Même énoncé si on met dans le coup un système local
d'anneaux sur $X'$. Enfin, $f$ induit \uncertain{une équivalence} entre la
catégorie des coéfficients locaux sur $C$ et celle des coéff.\ locaux sur
$C'$.
\note{la dactylographie a d'abord écrit « est pleinement fidèle (resp.\ une
équivalence) » puis effacé à la frappe « pleinement fidèle (resp. », de même
qu'à la ligne suivante « isomorphisme » et « foncteur pleinement fidèle
(resp.\ une équivalence) » se recouvrent ; on donne ce qui reste lisible. Les
lettres $C$, $C'$ de cette phrase sont tapées ainsi, là où le contexte attend
$X$, $X'$. Deux lignes entières sont ensuite effacées à la frappe.}
Principe de démonstration : on commence par prouver ce dernier résultat, en
notant que si un topos est localement hom.\ trivial, il est loc.\ connexe et
loc.\ simplement connexe, d'où une bonne théorie du $\pi_{0}$ et du $\pi_{1}$
(qui sont ici discrets), et on est ramené à un cas particulier du critère
d'homotopisme de Artin et Mazur, savoir un critère cohomologique pour qu'un
homomor\add{p}hisme de groupes $G \to H$ (ici de\add{s} groupes $\pi_{1}$ de
$X, X'$) soit un isomorphisme : il doit induire des isomorphismes sur les
$H^{0}$ et $H^{1}$ (y inclus dans le cas non commutatif…). On prouve la
pleine fidélité en se ramenant par la suite spectrale aux cas des
$\mathrm{Ext}^{i}$ de coéfficients locaux. Par la suite spectrale encore, cela
résultera du fait suivant : si $X$ est localement hom.\ trivial, alors la
catégorie des faisceaux abéliens loc.\ constants est stable par
$\underline{\mathrm{Ext}}^{i}$, et le foncteur $M \mapsto M_{X}$ (Ab) dans
$X_{\mathrm{ab}}$ commute aux dits $\mathrm{Ext}^{i}$.

En fait, $X$ et $X'$ étant loc.\ hom\add{.}\ triviaux, les conditions suivantes
sur $f$ seront équivalentes :
\begin{itemize}
\item[a)] $f$ est un homotopisme, i.e. induit pour tout système local (pas néc
commutatif) sur $X'$ un isomorphisme sur les $H^{i}$.
\item[b)] $f$ induit une équivalence
$D^{b}_{\mathrm{lc}}(X') \to D^{b}_{\mathrm{lc}}(X)$ .
\item[a')] $f$ induit une équivalence entre la catégorie des systèmes locaux
abéliens sur $X'$ et $X$, et des isomorphismes sur les $H^{i}$ correspondants
(do\add{nc} on ne prend ici que des coéfficients commutatifs).
\end{itemize}

\marginal{?}
J'ignore si on peut dans a) se borner aux systèmes locaux commutatif\add{s}.
L'équivalence entre a) et b) fournit une première justification ou motivation
pour définir des types d'homotopie via la catégorie
$D^{b}_{\mathrm{lc}}(X)$,
\note{le « $b$ » et le « lc » du dernier $D^{b}_{\mathrm{lc}}(X)$ sont écrits
à la main, en fin de ligne.}

\page{6}
éventuellement muni de la sous-catégorie pleine de tous les systèmes loc\add{aux}
sur $X$, et du foncteur cohomologique sur $D^{b}_{\mathrm{lc}}(X)$ à valeurs
dans ladite catégorie, et bien sûr du produit tensoriel (mais alors on sort de
$D^{b}$ pour entrer dans $D^{-}$ , redactor demerdetur).

\subsection*{II \quad $n$-catégories, catégories $n$-uples, et (Gr)-catégories.}

Appelons bi-catégorie une catégorie dans $(\mathrm{Cat})$, i.e. un triple
$(C_{0}, C_{1}, p_{1}, p_{2}, \pi)$, où $C^{\circ}$ est une catégorie,
$p_{1}, p_{2} : C_{1} \rightrightarrows C_{0}$ deux foncteurs,
$\pi : (C_{1}, p_{1}) \times_{C^{\circ}} (C_{1}, p_{2}) \to C_{1}$ un
foncteur, avec les axiomes habituels.
\note{les indices des $C_{i}$ et, plus bas, des $X_{ij}$ ont été tapés en
exposant ; il les biffe un à un à la main et les récrit en indice. On donne
partout sa correction ; les deux « $C^{\circ}$ » qu'il n'a pas touchés sont
laissés tels qu'ils sont tapés.}
Interprétant $(\mathrm{Cat})$ comme la sous-catégorie pleine de
$(\mathrm{Ssimpl})$ des objets qui … , et de même pour les catégories
à valeurs dans une catégorie telle $(\mathrm{Cat})$, on trouve que la
catégorie $(\mathrm{Bicat})$ des $(\mathrm{Cat})$-catégories est équivalente à
la sous-catégorie pleine de la catégorie des ensembles bi-semisimpliciaux
\[
\begin{array}{cccc}
X_{00} & X_{01} & \cdots & X_{0n} \\
X_{10} & X_{11} & \cdots & X_{1n} \\
\cdots & & & 
\end{array}
\]
formée des objets qui ligne par ligne et colonne par colonne satisfont la
condition d'ex\add{a}ctitude familière (transformer sommes amalgamées en
produit\add{s} fibrés). De même la catégorie $(\mathrm{Tricat})$ des
$(\mathrm{Bicat})$-catégories s'interprète en termes d'objets
tri-semi-simpliciaux, etc.

Exemple : soit $C$ une catégorie, posons $C_{0} = C$,
$C_{1} = \underline{\mathrm{Fl}}(C)$, avec structure\add{s} $p_{1}, p_{2}$,
évidentes (d'où une $(\mathrm{Cat})$-cat\add{ég}orie dont les composantes
semi-simpliciales supérieures sont les catégories
$\underline{\mathrm{Fl}}^{n}(C)$ des $n$-flèches de $C$ …) L'ensemble
bi-semisimplicial correspondant a pour composantes les ensembl\add{es}
$\mathrm{Hom}(\Delta^{n} \times \Delta^{m}, C)$. Par exemple, les objets de
$X_{11}$ sont les carrés commutatif\add{s} de flèches de $C$, qu'on peut
interpréter (de deux manières différentes) co\add{mme}

\page{7}
morphismes de flêches.

Notons une symmétrie évidente dans la catégorie $(\mathrm{Bicat})$,
interprété\add{e} en termes d'ensemble\add{s} bi-semisimpliciaux, obtenue en
échangeant les lign\add{es} et les colonnes dans de tels ensembles bi…
Un ensemble bisemisimpl\add{ici}al qui … peut donc de deux façons
différentes être considéré comme déf\add{i}nissant une bicatégorie. De même le
groupe symétrique $\mathrm{Sym}_{n}$ opère sur la catégorie
$(n\text{-upcat})$ des catégories $n$-uples ..

Une 2-catégorie peut être définie comme constituée par un ensemble $C_{0}$
(l'ensemble des objets ), qu'on interprètera comme une catégorie discrète,
plus la donnée pour tout $x, y \in C_{0}$ d'une catégorie
$\underline{\mathrm{Hom}}(x, y)$, plus des foncteurs composition entre
celles-ci … Prenant la catégorie somme $C$ des
$\underline{\mathrm{Hom}}(x, y)$, On peut exprimer la donnée de ces foncteurs
de composition et des sources et buts $x, y$ des catégories
$\underline{\mathrm{Hom}}(x, y)$, comme définis par des foncteurs « source »
et « but » $C_{1} \to C_{0}$, et un foncteur de composition
$C_{1} \times_{C_{0}} C_{1} \to C_{1}$ , les axiomes des 2-catégories
correspondant exactement au\add{x} axiomes des $(\mathrm{Cat})$-catégories. De
cette façon, la catégorie des 2-catégor\add{ies} est équivalente à la
catégorie des $(\mathrm{Cat})$-catégories qui ont une « catégo\add{rie}
d'objets » qui est discrète, ou encore à la catégorie des Bicatégories qui,
interprétées comme ensembles bi-semisimpliciaux, ont comme première ligne un
ensemble semisimplicial discret (i.e. un foncteur constant sur $\Delta$). Il
parait qu'on a une interprétation simple dans le même esprit de la notion de
$n$-catégorie en termes d'ensembles $n$-semi-simplici\add{aux}.
\note{Même correction qu'à la page précédente : les indices de $C_{0}$, $C_{1}$
tapés en exposant sont récrits à la main en indice.}

Demander à Quillen relations entre la notion de $m$-Catégorie dans la
catégorie $(n\text{-Cat})$, et celle d'ensemble $(m+n)$-semisimplicial.
\note{ce paragraphe est marqué d'un double trait vertical en marge de gauche.}

L'interprétation standard semisimpliciale permet d'interpréter

\page{8}
\marginal{Attention : la catégorie
$\underline{\mathrm{Hom}}_{\mathrm{grab}}$, qui s'identifie à la catégorie
définie par un $R\mathrm{Hom}_{\mathbb{Z}}(C_{x}, C'_{x})$, ce point devant
être élucidé !}
\note{note dans un cadre, en tête de la page, au-dessus du numéro 7 du
tapuscrit ; le signe qui l'ouvre la rattache sans doute à la dernière phrase
de la page précédente, mais aucun renvoi correspondant n'y est visible.
$C_{x}$, $C'_{x}$ ne sont pas définis ici.}

de même les $(\mathrm{Cat})$-groupes (ou groupes à valeurs dans
$(\mathrm{Cat})$) comme étant les groupes semisimpliciaux qui « sont » des
catégories i.e. qui satisfont la condition d'exactitude familière, permettant
d'en déte\add{r}miner tous les composants à l'aide de $C_{0}$ et $C_{1}$ .
Appliquant le tapis de Dold, dans le cas commutatif, on trouve que les
$(\mathrm{Cat})$-groupes abéliens forment une catégorie équivalente à celle des
complexes de chaines de groupes abéliens qui sont de longueur 1 :
\[
0 \leftarrow C_{0} \leftarrow C_{1} \leftarrow 0 \quad .
\]
Quand on explicite la $(\mathrm{Cat})$-groupe abélien associé à un tel complexe
de chaines, on trouve « ma » construction chérie de la « catégorie de Picard »
associée au complexe : objets ceux de $C_{0}$, flêches de $x$ dans $y$ les $z$
dans \add{$C_{1}$} tels que $dz = y - x$ . On trouve donc à isomorphisme près
tous les groupe\add{s} abéliens de $(\mathrm{Cat})$ ainsi. Un complément
important est que toute catégori\add{e} de Picard \struck{(resp.\ et
commutative)} est équivalente (au sens technique qu'on ne précis\add{e} pas
ici) à un vrai $(\mathrm{Cat})$-groupe \struck{(resp.\ $(\mathrm{Cat})$-groupe
commutatif)}. \add{Quand on \uncertain{simplifie} mod « \uncertain{Grab} »
équivalence, cela correspond à travailler dans la catégorie \emph{dérivée} des
complexes, seuls invariants \uncertain{sont} $H_{0}$, $H_{1}$
(\uncertain{ou} $\pi_{0}$, $\pi_{1}$).}
\marginal{\struck{\ill{}}}
\marginal{\ill{} si on part d'une « cat.\ de Picard \uncertain{commut.} »,
elle ne se \uncertain{réalise} \ill{} que par un $(\mathrm{Cat})$-g.\ \uncertain{non}
\uncertain{commut.}}
\note{l'insertion, en deux lignes serrées entre les lignes tapées, suit les
deux parenthèses biffées qui étendaient l'énoncé au cas commutatif. En regard,
dans la marge de gauche, deux blocs écrits en oblique : le premier, barré de
traits, ne se lit pas ; le second, pâle et rogné au bord gauche de la
photocopie, dit pourquoi les variantes commutatives sont biffées. L'exemplaire
annoté à l'encre (p.~38) porte les mêmes notes, plus lisibles.}
On notera qu'on a une classification beaucoup plus restrein\add{te} que pour les
$H$-espaces (resp.\ les $H$-espaces homotopiquement comm.) mod\add{ulo}
équivalence d'homotopie, de façon précise entre les groupes resp.\ les
groupes abéliens de la catégorie homotopique. La raison en est qu'on est
beaucoup plus restrictif dans le contexte catégorique pour le sens de
l'asso\add{cia}tivité (resp.\ et de la commutativité), puisqu'on ne demande pas
seulemen\add{t} l'existence d'isomorphismes d'associativité (qui suffirait
pour avoir u\add{n} groupe dans la catégorie homotopique
$(\mathrm{Cat})' \simeq (\mathrm{Ssimpl})'$, modulo l'existence d'unités et
d'invers\add{es}) mais qu'on suppose \emph{donnés} ces isomorphismes, de façon
à satisfaire à certaines compatibilités : c'est ce qui permet, me semble-t-il,
de prouver qu'\add{à} équivalence près, on peut supposer qu'il s'agit de vrais
groupes dans $(\mathrm{Cat})$, \struck{resp.\ de vraie commutativité.}
\note{« la catégorie homotopique » est tapé en interligne et mis en place par
un trait à la main ; une accolade et un trait en marge de gauche tiennent les
deux dernières lignes, dont la fin est biffée à la main.}

\page{9}
On doit trouver une approximation meilleure d'un cran à la classifi\add{ca}tion
topologique en prenant des $(2\text{-Cat})$-groupes resp.\ des
$(2\text{-Cat})$-group\add{es} abéliens, ou plus généralement des
« Gr-2-catégories » resp.\ « Grab-2-catégories » . (NB Une Grab-catégorie
définit une Grab-2catégor\add{ie} comme une catégorie définit une
2-catégorie ; du point de vue semi-simpli\add{ci}al, cette opération serait
analogue à celle de « cosquélette tronqué » ; Mais également une
Grab-2-catégorie définit une Grab-catégorie, savoir la catégorie des
automorphismes de l'objet unité ; cette opération serait analogue à celle de
« tuage des groupes d'homotopie en dim $> 1$ »). Pour l\add{a} notion
groissière (Gr ou Grab) l'idée de ma lettre à Verdier était qu\add{'on} doit
trouver la même chose que les tronqués au cran 2, dans la catégorie
\emph{homotopi\add{que}} des $H$-espaces --- i.e. ce qu'on déduit des
$H$-espaces par tuage des group\add{es} d'homotopie en dimensions $> 2$ .
\note{la frappe porte « définie en Grab-catégorie » ; « définit une » est la
correction écrite au-dessus, propre à cette copie (l'exemplaire annoté à
l'encre, p.~39, ne la porte pas). Les signes $>$ sont repassés à la main sur
une frappe surchargée, comme sur l'exemplaire à l'encre.}

(Demander à Quillen relations exactes entre objets-groupes et objet\add{s}
connexes de la catégorie homotopique ponctuée : forment-ils des catégories
équival\add{en}tes par le foncteur espace des lacets ?).
\marginal{non, on a un foncteur \uncertain{pl.} fidèle (?) mais pas
\uncertain{essentiellement} surjectif …}
\note{le paragraphe est tenu en marge par un double trait vertical ; la
réponse est écrite en oblique dans la marge de gauche.}

Pour les \emph{vrais} groupes abéliens de $(2\text{-Cat})$, l'approche de
Quillen consiste à les identifier d'abord à des groupes abéliens de la
catégorie $(\mathrm{Bicat})$, donc à des groupes abéliens bi-semisimpliciaux
particuliers (conditions sur chaque ligne et chaque colonne). Appliquant la
version bi-semisimpli\add{ci}ale du tapis de Dold, on trouve que la catégorie
de ces objets est équivalente à la catégorie des bicomplexes de chaines de
groupes abéli\add{ens} qui sont nuls en dehors du carré $0 \leq i, j \leq 1$.
La \struck{colonn\add{e}} \add{\uncertain{cond.}} que la première ligne corresponde
à un groupe semisimplicial trivial signifie alors qu'elle n'a que le terme de
degré zéro, i.e. que le bicomplexe est réduit à

\begin{tikzcd}
C_{00} & 0 \arrow[l] & 0 \arrow[l] \\
C_{10} \arrow[u] & C_{11} \arrow[l] \arrow[u] & 0 \arrow[l]
\end{tikzcd}

\note{le diagramme est complété à la main : les deux flèches verticales et
les « $0$ » de droite sont manuscrits. « cond. », au-dessus de « colonne »
biffé, est propre à cette copie : l'exemplaire à l'encre (p.~39) garde « La
colonne » tel quel.}

\page{10}
Ces bicomplexes s'identifient aux complexes de chaines de longueur 2. Donc la
catégorie de $(2\text{-Cat})$-groupes abéliens est équivalente à celle de\add{s}
complexes de chaines de longueur 2. L'explicitation de la 2-catégorie associée
à un tel complexe de chaines est mon interprétation chérie. Il faudrait encore
expliciter \add{la} 2-Grab-catégorie des homomorphismes (en un sens à
préciser) entre deux telles 2-Grab-caté\add{go}ries strictes, en termes d'un
complexe $R\mathrm{Hom}$.

L'idée qui se dégage de cette approche, et que Quillen on espère \add{va}
expliciter, c'est que la notion de groupe \emph{abélie\add{n}} dans
$(n\text{-Cat})$ correspond à celle de complexe de chaines de longueur $n$, en
raisonnant soit mod.\ isomorphisme de part et d'autre, soit mod.\ une notion
naturell\add{e} de $n$-équivalence d'un côté, et dans la catégorie dérivée de
l'autre. Cec\add{i} doit pouvoir se préciser en utilisant des
$\mathrm{Hom}^{\bullet}$ resp.\ des $R\mathrm{Hom}^{\bullet}$.

Si maintenant $C$ est une catégorie avec un « produit tensoriel »
assoc\add{ia}tif (et éventuellement unitaire), i.e. un bifoncteur avec des
isomorphismes d'associativité (mais sans inverse …), on peut lui associer
sans doute une Gr-catégorie universelle (tout comme à un monoïde on associe
\add{le} groupe symmétrisé), d'où des invariants, qu'on pourra désigner par
$K_{0}(C)$ et $K_{1}(C)$. On suppose qu'on pourra de même trouver, pour tout
$n$, une Gr-$n$-catégorie universelle où envoyer $C$ de façon compatible avec
$\otimes$ (tout cela dans un sens à préciser), dont les invariants $\pi_{n}$
seraient par déf\add{ini}tion les $K_{n}(C)$. Un cas intéressant est celui où
$C$ est une catégorie ave\add{c} produits finis (par exemple une catégorie
additive) et $\otimes$ est le produit, \add{[}Mais on notera que dans le cas d'une
catégorie non seulement additi\add{ve} mais abélienne, il y a une notion plus
adéquate du $K_{0}$ , qui devrait corr\add{es}pondre à une notion plus
adéquate de la suite des $K_{i}$ , inspirée du tapis de Roos ; sans doute
c'est en termes de (« vraie ») catégories triangulées
\note{le signe du produit tensoriel est obtenu à la machine en surchargeant
un O ; on le rend par $\otimes$. Le crochet ouvert avant « Mais » est tracé à
la main sur une parenthèse tapée ; il n'est pas refermé.}

\page{11}
qu'il conviendra de poser les définitions adéquates). --- Peut-être Quill\add{en}
a-t-il une définition semi-simpliciale plus directe, sans passer par des
$n$-catégories, des $K_{i}$ et plus précisément du $H$-espace qui devrait
correspo\add{n}dre à $C$. Ce n'est certainement pas en général l'ensemble
semi-simplicial $SC$ associé à $C$, puisque dans le cas d'une catégorie
additive, $C$ est homotopiquement trivial (en effet le foncteur identique est
« homoto\add{pe} » au foncteur nul, puisqu'il suffit d'avoir un homomorphisme
d'un foncteur dans l'autre ; d'où une équivalence d'homotopie entre $C$ et la
catégorie nulle). Pourtant, il semble bien que dans le cas où $C$ est une
Gr-catégor\add{ie} i.e. admet des inverses, alors ce soit bien la construction
« naïve » $SC$ qu\add{i} est la bonne, puisqu'on trouve bien un $\pi_{0}$ et un
$\pi_{1}$ qui sont aussi les $\pi_{0}$ et $\pi_{1}$ au sens catégorie sans
Gr-structure.
\marginal{oui}
\note{en marge de gauche, au droit de la question « Peut-être Quillen a-t-il
une définition semi-simpliciale plus directe », tenue par un double trait ; un
signe de renvoi, en face de « i.e. admet des inverses », appelle le
paragraphe manuscrit qui suit.}

\marginal{À expliquer : pour des Gr-$n$-catégories, les $\pi_{i}$ au sens
semi-simplicial sont aussi les $\pi_{i}$ au sens Gr-$n$-catégories, savoir les
objets à isom.\ près dans divers $\underline{\mathrm{Hom}}(1, 1)$ … .
\struck{Il y a aussi une interprétation de l'opération \uncertain{$\Omega$ :
\ill{} des lacets} en termes du foncteur d'opération du shift géométrique
évidente sur les $n$-catégories …}}
\note{paragraphe manuscrit, dans le blanc laissé au bas du feuillet tapé,
appelé par un signe en marge et tenu par un trait vertical ; sa seconde
moitié, à partir de « Il y a aussi », est barrée de traits obliques.}

\page{12}
\subsection*{III \quad Point de vue « motivique » en théorie du cobordisme.}

Soit $C$ la catégorie des variétés différentiables \marginal{(compactes ?)}
(pas nécessairem\add{ent} orientables), les morphismes étant les classes
d'homotopie d'applications continues. Si $B$ est une catégorie, on s'intéresse
aux couples $(F_{\bullet}, F^{\bullet})$ d'un foncteur covariant et d'un
foncteur contravariant de $C$ dans $B$, satisfaisant les conditions que pour
tout $X \in \mathrm{Ob}\, C$, on a $F_{\bullet}(X) = F^{\bullet}(X)$, et que si
on a un produit fibré ordinaire

\begin{tikzcd}
 & X' \arrow[dl, "g'"'] \arrow[dr, "f'"] & \\
X \arrow[dr, "f"'] & & Y' \arrow[dl, "g"] \\
 & Y &
\end{tikzcd}

avec $f$ et $g$ transversaux, alors on a $g^{\bullet} f_{\bullet} =
f'_{!}\, g'^{\bullet}$ (on pose $f_{\bullet} = F_{\bullet}(f)$,
$f^{\bullet} = F^{\bullet}(f)$) . La question est de trouver une $B$ et un
couple $(F_{\bullet}, F^{\bullet})$ qui soient « universels » dans u\add{n}
sens évident. Quillen prouve que la construction suivante donne une telle
solution.
\note{« (compactes ?) », écrit au-dessus de « différentiables », est propre
à cette copie : l'exemplaire annoté à l'encre (p.~43) ne le porte pas. « les
morphismes étant les classes d'homotopie d'applications continues » et « (on
pose …) » sont tapés en interligne et mis en place par un trait à la main. La
formule est tapée « $g^{\bullet} f_{\bullet} = f!\, g'^{\bullet}$ » ; on lit
$f'_{!}$ sans certitude sur l'accent.}

Pour toute variété $X \in \mathrm{ob}\, C$ , on définit un anneau commutatif
$B(X)$ (l'anneau de bordisme de $X$) en prenant comme éléments les
class\add{es} d'éléments de $C$ au dessus de $X$, pour la relation de
cobordisme : $f' : X' \to X$ et $f'' : X'' \to X$ sont équivalents, s'il
existe une variété à bord $Z$, et un $f : Z \to X$, tel que
$\text{bord } Z = X' \amalg X''$ et $f'$ et $f''$ sont induits par $f$. La
somme de $B(X)$ est la somme disjointe (c'est une loi de grou\add{pe} où tout
élément est d'ordre \add{2}), le produit est donné par le produit fibré, en
utilisant le théorème de transversalité de Thom pour bouger un peu $f''$, de
sorte à le rendre transversal à $f'$.
\note{le chiffre de « d'ordre 2 » est récrit à la main sur une frappe biffée.}
\marginal{NB L'anneau est de plus \emph{gradué} par la dimension relative
…}
Si on a $f : X \to Y$, on définit $f_{\bullet} : B(X) \to B(Y)$ de façon
évidente (tou\add{te} variété sur $X$ est sur $Y$), et
$f^{\bullet} : B(Y) \to B(X)$ par produit fibré (util\add{i}sant encore le
th.\ de transversalité). De cette façon, $B(X)$ devient u\add{n}
foncteur-anneau \add{(gradué)} contravariant en $X$, et un foncteur
groupe covariant e\add{n}

\page{13}
$X$. De plus, on a la formule \add{(du type
$g^{\bullet} f_{\bullet} = f'_{!}\, g'^{\bullet}$ ci-dessus, impliquant la
formule)} de projection habituelle. Ceci posé, les constructions bien connues en
théorie des correspondances permettent de définir une catégorie
$\underline{B}$ dont les objets sont ceux de $C$, et les $\mathrm{Hom}$ étant
donnés par
\[
\mathrm{Hom}_{\underline{B}}(X, Y) = B(X \times Y) \quad .
\]
On a un isomorphisme canonique $\underline{B} \simeq \underline{B}^{\circ}$ ,
qui est l'identité sur les objets et la symmétrie
$B(X \times Y) \simeq B(Y \times X)$ sur les flêches, et un foncteur évident
$F_{\bullet} : C \to \underline{B}$ , qui est l'identité sur les objets et
don\add{né} par les graphes sur les flêches. Composé avec la symmétrie
précédente on trouve un foncteur $F^{\bullet} : C \to \underline{B}$. On
vérifie enfin que le couple $(F_{\bullet}, F^{\bullet})$ satisfait aux
conditions envisagées dans l'alinéa 1 (ceci aus\add{si} devrait être formalisé
au cas des correspondances générales …). On vérifie enfin qu'il a la
propriété universelle qu'il faut. Lorsqu'on se donne un couple
$(F_{!}, F'^{\bullet})$ à valeurs dans une catégorie $B'$, il faut trouver une
factorisation par $B$. Elle est toute trouvée sur les objets, il faut la
définir sur les flêches, donc définir, pour $X, Y$ dans $C$, une application
\[
(*) \qquad B(X \times Y) \to \mathrm{Hom}(X', Y') \quad .
\]
Mais si $Z$ est au dessus de $X \times Y$, d'où $p : Z \to X$ et
$q : Z \to Y$, on lui assoc\add{ie} un élément du deuxième membre de $(*)$
comme le composé
\[
X' \xrightarrow{\;F'^{\bullet}(p)\;} Z' \xrightarrow{\;F_{!}(q)\;} Y'
\quad .
\]
Il faut prouver que l'élément obtenu en remplaçant $Z$
\add{$(= Z_{0})$} par un $Z_{1}$ cobordant au dessus de $X$ \add{est le
même}. Or si $g : T \to X \times Y$ réalise le cobordisme, on voit (regardant
seulement $T$ et la décomposition $\mathrm{bord}(T) = Z \amalg Z_{1}$) qu'il
existe une variété $\overline{T}$ sur $S^{1}$ (obtenue en « doublant » $T$)
transversale au dessus d'un couple de points $s_{0}, s_{1}$ de $S^{1}$, tels
que $Z_{i}$ soit l'image inverse \add{dans $\overline{T}$} de $s_{i}$
($i = 0, 1$) ; d'autre part,
\note{la parenthèse manuscrite en tête de la page est rattachée par un trait
à « la formule ». « $(= Z_{0})$ », au-dessus de $Z$, est pâle sur la
photocopie ; « est le même », écrit sous « l'élément obtenu », est renvoyé
par un trait après « au dessus de $X$ ».}

\page{14}
$g$ manifestement se prolonge en $\overline{g} : \overline{T} \to X \times Y$ .
Cela nous ramène à prouv\add{er} que le composé
$\overline{T}' \xrightarrow{\;\alpha_{i}\;} Z'_{i}
\xrightarrow{\;\alpha_{i*}\;} \overline{T}'$ ne dépend pas de
$i = 0, 1$. Ecrivons $X$ au lieu de $\overline{T}'$, considérons le diagramme
cartésien d'inclusions

\begin{tikzcd}
X_{s} \arrow[r, "\alpha_{s}"] \arrow[d, "\alpha_{s}"'] & X \arrow[d, "{\Gamma = (p, \mathrm{id}_{X})}"] \\
X \arrow[r, "i_{s}"'] & S^{1} \times X
\end{tikzcd}

on a alors $\alpha_{s}^{*}\, \alpha_{s*} = \Gamma^{*}\, i_{s*}$ , qui ne
dépend pas de $s$ puisque $i_{s}$ n'en dépend pas à homotopie près. --- Le
reste est sans doute l'AQT.
\note{les étiquettes du composé, les flèches et les $\alpha_{s}$ du carré,
l'étiquette $\Gamma = (p, \mathrm{id}_{X})$ de la flèche de droite et les
astérisques de la formule sont portés à la main dans les blancs laissés par la
machine ; petits et pâles sur la photocopie, ils se lisent comme sur
l'exemplaire annoté à l'encre (p.~45). La première étiquette du composé est
$\alpha_{i}$, sans astérisque visible.}

Sur la str\add{u}cture des $B_{\bullet}(X)$. Ce sont tous des algèbres graduées
sur $B_{\bullet}(\mathrm{pt}) = \Lambda$ , qui est une algèbre graduée à
degrés positifs, réduit\add{e} à $\underline{\mathbb{F}}_{2}$ en degré 0. C'est
d'ailleurs une algèbre de polynômes \add{(à une infinité de générateurs
$x_{i}$ ($i$ \uncertain{entier $\geq 1$}, \uncertain{$\neq$ de la forme
$2^{\nu} - 1$}))} connue …
\note{l'insertion, appelée par un signe en marge, est écrite en deux lignes
sous la ligne tapée ; la seconde, pâle et rognée au bord droit de la
photocopie, se lit sur l'exemplaire à l'encre (p.~45).}

On a un homomorphisme évident d'anneaux gradués $B_{\bullet}(X) \to H^{\bullet}(X)$,
nul en degrés $> 0$, qui à tout $Z$ sur $X$ associe l'image de
$1 \in H^{0}(Z)$ par l'homomorphisme de Gys\add{in} associé à $Z \to X$. Cet
homomorphisme par Thom est surjectif (toute cla\add{s}se de cohomologie
provient d'une image de variété compacte …) ;
\marginal{NB $H^{*}(X)$ ($= H^{*}(X, \underline{\mathbb{F}}_{2})$)}
\marginal{en fait $H^{\bullet}(X) \simeq B_{\bullet}(X) / \Lambda^{+}
B_{\bullet}(X)$}
\note{deux notes superposées dans la marge de gauche : « NB $H^{*}(X)$ … »
est propre à cette copie ; « en fait … » est la note de l'exemplaire à l'encre
(p.~45). « d'anneaux gradués » est tapé en interligne ; l'exposant de $H$ est
une frappe surchargée.}

On prouve même qu'il existe un relèvement fonctoriel en $X$
$H^{\bullet}(X) \to B_{\bullet}(X)$, \add{(il est douteux qu'il soit)}
compatible avec multiplication, tel que l'on ait :
\note{la main entoure « compatible avec multiplication » d'une parenthèse et
écrit au-dessus « il est douteux qu'il soit » ; un point d'interrogation dans
la marge.}

\struck{a)} L'homomorphisme correspondant
$H^{\bullet}(X) \otimes_{\underline{\mathbb{F}}_{2}} \Lambda \simeq
B_{\bullet}(X)$ est un isomorphisme (ce sera vrai pour tout relèvement
individuellement …), e\add{t}

\struck{b)} \add{[Quillen ignore si} \struck{Le} relèvement est compatible
aussi avec les hom.\ de Gysin
\note{les lettres a) et b) sont biffées à la main ; le crochet ouvert à
l'encre devant b) fait de la seconde clause une question, non plus une
condition. Le bord droit est rogné : la fin de « et » et le crochet fermant
que porte l'exemplaire à l'encre (p.~45) n'y paraissent pas. Un grand point
d'interrogation dans la marge, en regard.}

Ainsi, la catégorie $\underline{B}$ admet un foncteur pleinement fidèle dans la
catégorie des modules gradués libres de type fini sur $\Lambda$, les modules de
type fini qu'on obtient ainsi ayant des générateurs à degrés $< 0$ satisfaisant
(pour $X$ connexe i.e. quand il y a un seul générateur de degré zéro) la
dualité de Poincaré. Quand on prend la catégorie \emph{quasi abélienne}
\add{$\widetilde{\underline{B}}$} associée à $\underline{B}$, on trouve donc
tous les modules gradués libre\add{s} de type fini à générateurs en degrés
$\leq 0$ seulement. C'est u\add{ne}
\marginal{i.e.\ l'enveloppe « karoubienne »}
\note{en oblique dans la marge de gauche, en regard de « quasi abélienne » ;
cette note est propre à cette copie.}

\page{15}
catégorie additive (et même quasi-abélienne) avec $\otimes$ qui a toutes les
vertus, sauf d'être abélienne, et sauf d'admettre des
$\underline{\mathrm{Hom}}$ internes.
\marginal{quasi-abélienne $=$ karoubienne $+$ additive}
\note{« additive (et même quasi-abélienne) » est tapé en interligne ; la note
qui le commente, écrite au-dessus en tête de la page, est propre à cette
copie.}
Pour trouver ces derniers, on peu\add{t} si on veut paraphraser la
construction de la catégorie des motifs, \add{en} considérant l'objet
\add{$\Lambda(-1)$} de $\widetilde{\underline{B}}$ conoyau de
$p^{*} : e \to S$, où $p : S^{1} \to e$ est l'application du cercle sur un
point. C'est un objet de poids 1 (le poids est défini par le degré des
générateurs, changés de signe). On constate aussitôt que le foncteur
$M \mapsto M(-1) = M \otimes \Lambda(-1)$ de $\widetilde{\underline{B}}$ dans
lui-même est pleinement fidèl\add{e} ce qui permet alors, par une construction
formelle standard, de com\add{plé}ter $\widetilde{\underline{B}}$ en une
catégorie où $\Lambda(-1)$ devient inversible. Il se trouve alors que dans
cette dernière les $\underline{\mathrm{Hom}}(X, Y)$ internes existent et sont
isomorphes à $\check{X} \otimes Y$, … en fait , on trouve maintenant \add{la}
catégorie de \emph{tous} les modules gradués libres de type fini sur
\add{$\Lambda$}. Bien sûr, cette catégorie n'est pas encore abélienne pour
autant, puisque \add{$\Lambda$} n'est pas un corps ; si on veut définir la
catégorie abélienne engendrée, on s'attend à trouver la catégorie des modules
de \struck{type} \add{présentation} f\add{inie} sur $\Lambda$. (\add{qui est
\uncertain{abélienne} bien que $\Lambda$ ne soit pas noethérien, car $\Lambda$
est} \struck{est-il noethérien, i.e. le nombre de générateurs est-il fini ?})
\add{\uncertain{lim} d'anneaux noeth.\ à homs de transition \emph{plats}}).
\note{« si on veut » est tapé en interligne ; « $\Lambda(-1)$ » et les tildes
sur $\underline{B}$ sont portés à la main. Trois interventions sont propres à
cette copie : « en », écrit au-dessus de « considérant » là où la frappe
« en », en fin de ligne, est rognée, et les deux $\Lambda$ de « sur
$\Lambda$ » et « puisque $\Lambda$ n'est pas un corps », tracés à la main dans
des blancs que la machine avait laissés et que l'exemplaire à l'encre (p.~46)
laisse vides. La question tapée sur la noethérianité de $\Lambda$ est biffée
et la réponse écrite au-dessus et à la suite, signalée d'un signe en marge.}

Cette construction ressemblerait à celle d'une théorie des motifs où on
n'aurait pas passé à l'équivalence homologique pour les cycles, mais où on
aurait gardé l'équivalence rationelle, --- qui est bien trop fine à beaucoup
d'égards. Pour trouver un équivalent plus étroit, il \uncertain{convient} de
reprendre la théorie précédente, on pourrait formuler le même Pb universel que
plus haut, mais en prena\add{nt} l'équivalence homologique (via le graphe) des
morphismes de variété\add{s} au lieu de l'équivalence homotopique.
\struck{Mais alors il semble que l'on trouve comme catégorie solution celle
dont les objets sont les vari}

\page{16}
\struck{tés, les morphismes \emph{toutes} les correspondances cohomologiques (en vertu du
théorème de Thom rappelé plus haut). On trouve donc un foncteur pleinement
fidèle de cette catégorie dans celle des vectoriels gradués positivement de
dim.\ finie sur $\underline{\mathbb{F}}_{2}$ , qui est déjà abélienne, et si
on veut rendre comme plus haut $\underline{\mathbb{F}}_{2}(-1)$ inversible, on
trou\add{ve} tous les vectoriels gradués de dim.\ finie sur
$\underline{\mathbb{F}}_{2}$ . C'est un peu tro\add{p} trivial comme structure,
et on peut faire mieux,}
\note{ces sept lignes, qui achèvent le passage biffé au bas de la page
précédente, sont barrées à la main de grandes croix ; « co » y est tapé
au-dessus de « homologiques ».}

Pour trouver une catégorie universelle, on associe à toute $X$ $S(X)$
l'ensemble des sous-variétés de $X$, modulo équivalence homologique. Donc
$S(X) \subset H^{\bullet}(X)$, et on prouve qu'on trouve en fait un sous
anneau de $H^{\bullet}(X)$ (grâce à un théorème de Thom qui caractérise $S(X)$
e\add{n} termes d'opérations de Steenrod, cf plus bas). On prouve de même (par
le même théorème) que les $S(X \times Y)$ sont stables par composition
\add{des} correspondances cohomologiques (cela ne semble pas trivial, puisque
les $S(X)$ ne sont pas stables par Gysin…). D'où une catégorie
add\add{i}tive dont les objets sont ceux de $C$, les flêches sont les
\add{\uncertain{él.\ des}} $S(X \times Y)$. Il semble qu'elle doive jouer le rôle
universel que j'ai envisagé plus haut. En tous cas, sa définition géométrique
est fort simple\add{.} Nous allons en donner une interp\add{r}étation
algébrique. Pour ceci, considérons l'anneau des opérations cohomologiques
« stables » \struck{des endomorphismes du foncteur $X \mapsto H^{\bullet}(X)$
qui sont additifs, et qui en degré zéro se réduisent à un multiple de
l'identité}. C'est un anneau gradué, avec application diagonale
anticommutative i.e. commutative (car.\ 2 !), associative et unitaire. So\add{n}
dual gradué est une algèbre graduée anticommutative donc commutative, avec
application diagonale associative et unitaire ; il est réduit à
$\underline{\mathbb{F}}_{2}$ en degré zéro. Donc il correspond à un schéma en
groupes affine $G$ sur $\underline{\mathbb{F}}_{2}$ , qui est
d'aille\add{urs} le foncteur des $\otimes$-automorphismes additifs du foncteur
$H^{\bullet}$ sur $C$ \add{qui opèrent trivialement sur \struck{\ill{}}
\uncertain{$H^{1}(S^{1})$}}. \struck{\ill{}}
\note{« $S(X)$ », après « à toute $X$ », est tapé en interligne et mis en
place par un trait, la frappe « $S(X)$ » après « sous-variétés » étant biffée à
la machine. « él.\ des », au-dessus de « $S(X \times Y)$ », est propre à
cette copie : l'exemplaire à l'encre (p.~47) ne le porte pas. « stables »
remplace à la machine un mot biffé ; la définition tapée qui suit est biffée
et repassée à la main d'un trait sinueux, avec un signe en marge. « i.e.\ commutative
(car.\ 2 !) », « graduée anticommutative donc » et « additifs »
sont tapés en interligne et mis en place par des traits à la main. La phrase
est achevée à la main au bas de la page ; l'exposant de $H(S^{1})$ est récrit
et se lit $1$ sans certitude.}

\page{17}
Le théorème de Thom est alors qu'un $x \in H^{\bullet}(X)$ est dans $S(X)$ sss
il est fixe par $G$, i.e. satisfait $A^{+} . x = 0$. (Pourquoi cette condition
est-el\add{le} déjà nécessaire ?) Une autre façon de le voir est de dire que le
foncteur $H^{\bullet}$ considéré comme allant de la catégorie
$\underline{S}$ précédente dans la catégorie des représentations de $G$ dans
des vectoriels sur $\underline{\mathbb{F}}_{2}$ , est pleinement fidèle. Quand
on prend l'enveloppe pseudo-abélienne, on trouve une sous-catégorie pleine de
la catégorie des représentations de $G$, mais il est douteux qu'on trouve une
catégorie abélienne, voire la catégorie de toutes les représentations de $G$
de dimensi\add{on} finie sur $\underline{\mathbb{F}}_{2}$ . (Quelle est la
structure de $G$, qui est prolisse ; d'après Quillen, il serait
pro-unipotent …). Ici $G$ bien sûr joue le rôle d'u\add{n} groupe de
Galois motivique, mais contrairement à la situation familière, il semble qu'il
soit fortement non réductif (et d'ailleurs le corps de définition dans tout
ceci est $\underline{\mathbb{F}}_{2}$ , et non $\underline{\mathbb{Q}}$ ! ).
\note{« fixe » est tapé en interligne au-dessus d'un mot biffé à la machine.}
\marginal{oui, car son anneau affine est une \uncertain{algèbre de polynômes}
à une infinité d'indéterminées, \struck{\ill{}} \uncertain{réunion
d'algèbres de polynômes} à \ill{} \ill{} \uncertain{finis} stables par
\uncertain{$\Delta$}}
\note{bloc de lignes courtes écrites en oblique dans la marge de gauche, en
regard de « d'après Quillen, il serait pro-unipotent » et des lignes qui
suivent, et relié au texte par un trait ; pâle et rogné au bord gauche de la
photocopie, il se lit sur l'exemplaire à l'encre (p.~48), dont la fin reste
incertaine elle aussi.}

Quillen\struck{t} donne une généralisation de la construction d'une catégorie
universelle $\underline{B}$ au cas où on travaille avec des variétés pas
nécessa\add{i}rement propres (mais $f_{\bullet}$ n'étant supposé défini que
pour $f$ propre), et quand on introduit une condition d'orientation
stable sur les \emph{morphismes} envisagés , via un morphisme
$H \to BO(\infty)$, où $H$ est un $H$-espace (commutatif, associatif et
unitaire ?) : une orientation de $f : X \to Y$ sera un relèvement homotopique
de $X \to BO$, défini par $T_{X} - f^{\bullet}(T_{Y})$ , en $X \to H$ les
éléments de $H(X, Y)$ seront définis par des diagrammes de morphismes orientés
$Z \to X$, $Z \to Y$, le deuxième étant propre, modulo bords de tels. En tant
qu'ensemble, il se définit aussi en termes de la variété $X \times Y$ munie
simplement de la famille des supports formée des ensembles fermés propres sur
$Y$ … .
\note{la frappe porte « Quillent » ; le « t » est barré à la main sur cette
copie seulement. « stable » est tapé en interligne.}

\marginal{Les $B^{H}_{m}(X)$ ($X$ compact \uncertain{disons}) peuvent se
décrire homotopiquement à la Thom, comme les groupes d'homotopie stables
d'espaces $\underline{\mathrm{Hom}}(X, MH(N - m))$ -----.}
\note{à la main, à la suite du texte tapé, au pied de la page.}

\marginal{cf exposé Karoubi à Bourbaki !}
\note{au bas de la page, soulignée ; cette note est propre à cette copie.}

\end{document}
